% Problem record op_02bb8f8228649ac3. This TeX file is derived from
% database/problems_json/op_02bb8f8228649ac3.json by scripts/sync-tex.mjs; edit the JSON record
% and rerun the script rather than editing this file.
\section{Advantage of fully general superchannel-discrimination strategies}
\label{sec:op_02bb8f8228649ac3}

\paragraph{Problem.}
Can a fully general adaptive strategy attain a larger Stein exponent than
every nested-adaptive strategy for discriminating two quantum superchannels?
Fix finite-dimensional physical realizations
\begin{equation}
  \Theta_i(\mathcal N)
  =\mathcal D_i\circ(\mathcal N\otimes\operatorname{id}_{S_i})
     \circ\mathcal E_i,
  \qquad i\in\{1,2\}.
  \label{eq:p67-superchannel-realizations}
\end{equation}
The systems $S_i$ in Eq.~\eqref{eq:p67-superchannel-realizations} are internal
memories.  A nested strategy recursively places complete uses of $\Theta_i$ inside one
another.  A fully general strategy may interleave the preprocessing and
postprocessing components of different uses in any causally valid order,
provided each $\mathcal E_i$ precedes its matched $\mathcal D_i$ and the
tester cannot access the internal memory $S_i$.  Define the vanishing-type-I-error Stein
exponent for a strategy class $\mathsf S$ by
\begin{equation}
  \zeta_{\mathsf S}(\Theta_1\|\Theta_2)
  :=\lim_{\varepsilon\downarrow0}\liminf_{n\to\infty}
    -\frac1n\log_2
    \inf_{\substack{P\in\mathsf S_n:\alpha_n(P)\leq\varepsilon}}
    \beta_n(P).
  \label{eq:p67-strategy-exponent}
\end{equation}
The definition in Eq.~\eqref{eq:p67-strategy-exponent} uses the type-I and
type-II errors $\alpha_n(P)$ and $\beta_n(P)$ of protocol $P$.
Does there exist a pair of realizations for which
\begin{equation}
  \zeta_{\mathrm{fg}}(\Theta_1\|\Theta_2)
  >\zeta_{\mathrm{nest}}(\Theta_1\|\Theta_2),
  \label{eq:p67-general-advantage}
\end{equation}
where $\mathrm{fg}$ denotes fully general strategies?  If not, prove equality
in Eq.~\eqref{eq:p67-general-advantage} with $>$ replaced by $=$ for all
superchannel pairs.

\subsection*{Status}

Unsolved

\subsection*{Source}

Hirche explicitly identified the achievability and optimality of fully general
adaptive strategies as the principal unresolved problem in quantum-network
discrimination \sourcecite{ref:p67-hirche}{Hir23}.

\subsection*{Progress}

\begin{itemize}
  \item Fully parallel strategies are exactly characterized by the regularized
  superchannel relative entropy.  Since they are contained in the fully
  general class, Hirche's meta-converse yields
  \begin{equation}
    D_{\rm sc}^{\infty}(\Theta_1\|\Theta_2)
    \leq\zeta_{\mathrm{fg}}(\Theta_1\|\Theta_2)
    \leq D_{\rm sc}^{A*}(\Theta_1\|\Theta_2),
    \label{eq:p67-general-bounds}
  \end{equation}
  where $D_{\rm sc}^{A*}$ is a fully amortized superchannel relative entropy
  that accounts for the exposed preprocessing and postprocessing components
  \sourcecite{ref:p67-hirche}{Hir23}.

  \item A braided strategy gives a concrete fully general ordering that need
  not be nested.  No pair is known for which such an ordering makes the first
  inequality in Eq.~\eqref{eq:p67-general-bounds} strict, and no protocol is
  known to attain $D_{\rm sc}^{A*}$ in general
  \sourcecite{ref:p67-hirche}{Hir23}.

  \item For classical superchannels, the fully amortized relative entropy
  collapses to the ordinary superchannel relative entropy.  Product strategies
  are therefore optimal, ruling out Eq.~\eqref{eq:p67-general-advantage} in
  the classical special case \sourcecite{ref:p67-hirche}{Hir23}.
\end{itemize}

\subsection*{References}

\begin{enumerate}[label={},leftmargin=4.8em,labelsep=0.5em]
  \footnotesize
  \item[\textup{[Hir23]}]\label{ref:p67-hirche}
  C. Hirche, ``Quantum Network Discrimination,''
  \emph{Quantum} \textbf{7}, 1064 (2023).
  \href{https://doi.org/10.22331/q-2023-07-25-1064}{doi:10.22331/q-2023-07-25-1064};
  \href{https://arxiv.org/abs/2103.02404}{arXiv:2103.02404}.
\end{enumerate}

\subsection*{Comment}

The open alternatives are operationally different: either construct a strict
advantage as in Eq.~\eqref{eq:p67-general-advantage}, prove a protocol that
attains the upper bound in Eq.~\eqref{eq:p67-general-bounds}, or sharpen that
converse until it meets an achievable nested or parallel rate.

\subsection*{Field}

Quantum metrology

\subsection*{Topic}

Channel discrimination; Quantum hypothesis testing; Superchannels and quantum combs

\subsection*{ID}

\texttt{op\_02bb8f8228649ac3}
